\documentclass{beamer}

\usepackage[utf8]{inputenc}
\usepackage[T1]{fontenc}
\usepackage{textpos}
\usepackage{tikz}

\usetheme{Madrid}
\usecolortheme{beaver}

% Custom changes:
\setbeamertemplate{footline}[frame number]{}
\definecolor{university_tuebingen}{RGB}{165,30,55}
\setbeamercolor{frametitle}{fg=university_tuebingen, bg=white}
\setbeamercolor{title}{fg=university_tuebingen}

\addtobeamertemplate{frametitle}{}{
\begin{tikzpicture}[remember picture, overlay]
\node[anchor=north east,yshift=0cm] at (current page.north east)
{\includegraphics[width=3cm]{../common/logo_uni_tuebingen2.png}};
\end{tikzpicture}}



\author{Prof. Dr. Christiane Zarfl, Dipl.-Inf. Willi Kappler}
%\date{\today}
\date{}
%\institute{Universität Tübingen}
\institute{}


\matlabTitle{4. Schleifen}

\setcounter{mchapter}{4}
\setcounter{mexercise}{0}

\begin{document}
  {
\beamertemplatenavigationsymbolsempty % suppress navigation on this (= first) slide
\begin{frame}[plain] % plain means: no header and footer on this (= first) slide
    \begin{textblock*}{0cm}(0.5cm, -0.7cm)
        \includegraphics[width=11.0cm]{../common/logo_uni_tuebingen.png}
    \end{textblock*}
    \titlepage
    \begin{textblock*}{0cm}(0.1cm, -2.5cm)
        \textcolor{university_tuebingen}{\rule{11.8cm}{0.2cm}}
    \end{textblock*}
\end{frame}
}


  \section{Einleitung}

  \subsection{Motivation}
  \begin{frame}
      \frametitle{Sie wissen bereits...}
      \begin{itemize}
          \item wie Sie durch Scripte Befehlsfolgen wiederverwertbar machen.
      \end{itemize}

      \textit{Wie kann ich häufig vorkommende Berechnungen/Abläufe automatisieren?}
  \end{frame}

  \begin{frame}
      \frametitle{Nach diesem vierten Block...}
      \begin{itemize}
          \item können Sie Berechnungen ``umgangssprachlich'' als Algorithmus formulieren.
          \item können Sie in Matlab Algorithmen implementieren und verwenden dabei sicher die Hilfsmittel der
          \begin{itemize}
              \item logischen Operatoren
              \item Wenn-Dann-Anweisungen
              \item While- und For-Schleifen
          \end{itemize}
      \end{itemize}
  \end{frame}

  \section{Programmieren}

  \subsection{Vorgehen}
  \begin{frame}
      \frametitle{Matlab als Programmiersprache}
      \begin{itemize}
          \item \textbf{Ziel}: Lösen einer (beliebig) komplizierten Berechnungsaufgabe
          \item \textbf{Vorgehen}:
          \begin{enumerate}
              \item Genaue Beschreibung der Aufgabe
              \item Formulierung eines \textbf{Algorithmus} (Schritt-für-Schritt Berechnungsanweisung) in Form eines Ablaufplans/``Kochrezepts'' oder Pseudocode
              \item Umsetzung in Programmiersprache
              \item Testen des Programms, evtl. zurück zu Punkt 2, oder 3
              \item Wartung und Pflege des Programms während der Nutzung
          \end{enumerate}
          \item Setzt voraus, dass Möglichkeiten der Programmiersprache bekannt sind
          \item In diesem Kurs geht es hauptsächlich um Punkt 3
      \end{itemize}
  \end{frame}


  \subsection{Pseudocode}
  \begin{frame}
      \frametitle{Pseudocode}
      \begin{itemize}
          \item Sprachliche Mischung aus natürlicher Sprache, mathematischer Notation und einer höheren Programmiersprache
          \item Dient genauer Beschreibung des Algorithmus
          \item Ist für Menschen leicht verständlich
          \begin{itemize}
              \item kann aber vom Computer (noch) nicht ausgeführt werden
              \item Ist inde den meisten Fällen keine Programmiersprache
              \item Orientiert sich aber oft an ``echten'' Programmiersprachen
          \end{itemize}
          \item Soll Algorithmen verständlich und klar ausdrücken, ohne auf die Eigenheiten einer Programmiersprache Rücksicht nehmen zu müssen
          \item Umgangssprache hilft Verfahrensschritte zu verdeutlichen
          \begin{itemize}
            \item \texttt{``durchlaufe das Feld a mit Index i''}
            \item \texttt{``vertausche die Inhalte der Variablen x und y''}
          \end{itemize}
      \end{itemize}
  \end{frame}

  \subsection{Bsp.: Telefonieren}
  \begin{frame}
      \frametitle{Bsp.: Telefonieren}
      \begin{itemize}
          \item Telefoniervorgang als Pseudocode:
          \lstinputlisting[language=pseudo]{pseudocode.txt}%
      \end{itemize}
    \end{frame}

    \subsection{Flussdiagramm}
    \begin{frame}
        \frametitle{Flussdiagramm}
        \begin{itemize}
          \item Telefoniervorgang als Flussdiagramm:
        \end{itemize}
        \begin{tikzpicture}[node distance = 2.5cm]
          \node [block] (init) {Hörer abnehmen};
          \node [decision, right of=init] (free) {Frei-zeichen};
          \node [block, right of=free] (dial) {wähle};
          \node [decision, right of=dial] (long) {langes Tuten};
          \node [block, right of=long] (wait) {auf Antwort warten};
          \node [decision, below of=long] (short) {kurzes Tuten};
          \node [block, right of=short] (hangup) {lege auf};
          \node [block, below of=free, text width=5em] (check) {lege auf und überprüfe Anschluss};
          \path [line] (init) -- (free);
          \path [line] (free) -- node [above] {ja} (dial);
          \path [line] (free) -- node [left] {nein} (check);
          \path [line] (dial) -- (long);
          \path [line] (long) -- node [above] {ja} (wait);
          \path [line] (long) -- node [left] {nein} (short);
          \path [line] (short) -- node [above] {ja} (hangup);
        \end{tikzpicture}
      \end{frame}

      \secMexercise
      \begin{frame}
          \frameMexercise
          \begin{exercise}
              \sloppy
              Schreiben Sie einen Pseudocode für die Berechnung des Mittelwerts und der Standardabweichung \textit{s} eines Zahlenvektors \textbf{x} \\
              \begin{displaymath}
                m = \frac{1}{n} \sum_{i=1}^{n}X_{i}
              \end{displaymath}
              \begin{displaymath}
                s = \sqrt{\frac{1}{n-1} \sum_{i=1}^{n}\left( x_{i} - m \right)^{2}} = \sqrt{\frac{1}{n-1} \left( \sum_{i=1}^{n} x_{i}^{2} - nm^{2} \right)}
              \end{displaymath}
          \end{exercise}
      \end{frame}

      \subsection{Logische Operationen}
      \begin{frame}
          \frametitle{Logische Operationen}
          \begin{itemize}
            \item Vergleich zwischen zwei \textbf{Zahlen}:
            \begin{itemize}
              \item \matlabInput{a=(3>1)} $\Rightarrow$ erzeugt logische Variable \texttt{a} mit Wert \matlabOutput{true (1)}
              \item \matlabInput{2>=4} $\Rightarrow$ \matlabOutput{false (0)}
              \item \matlabInput{1<pi} $\Rightarrow$ \matlabOutput{true (1)}
              \item \matlabInput{3==pi} $\Rightarrow$ \matlabOutput{false (0)}
              \item vergleichendes ``ist gleich'' durch \texttt{``==''}, Zuweisung durch \texttt{``=''}
            \end{itemize}
            \item Vergleich zwischen zwei \textbf{Vektoren}:
            \begin{itemize}
              \item \matlabInput{[1:4]<[5:-2:-1]} $\Rightarrow$ \matlabOutput{[true,true,false,false]}
            \end{itemize}
            \item Logisches ``und'' durch ``\string&'':
            \begin{itemize}
              \item \matlabInput{(x>1)&(x<5)} $\Rightarrow$ wahr für alle Werte von \texttt{x} zwischen 1 und 5
            \end{itemize}
            \item Logisches ``oder'' durch ``\string|'':
            \begin{itemize}
              \item \matlabInput{(x>1)|(x<5)} $\Rightarrow$ ist immer wahr
            \end{itemize}
            \item Logische Negation durch ``\string~'':
            \begin{itemize}
              \item \matlabInput{a~=5} $\Rightarrow$ wahr, wenn a $\neq$ 5
            \end{itemize}
          \end{itemize}
      \end{frame}


\end{document}
